% Options for packages loaded elsewhere
\PassOptionsToPackage{unicode}{hyperref}
\PassOptionsToPackage{hyphens}{url}
\PassOptionsToPackage{dvipsnames,svgnames*,x11names*}{xcolor}
%
\documentclass[
  11pt,
]{article}
\usepackage{lmodern}
\usepackage{amssymb,amsmath}
\usepackage{ifxetex,ifluatex}
\ifnum 0\ifxetex 1\fi\ifluatex 1\fi=0 % if pdftex
  \usepackage[T1]{fontenc}
  \usepackage[utf8]{inputenc}
  \usepackage{textcomp} % provide euro and other symbols
\else % if luatex or xetex
  \usepackage{unicode-math}
  \defaultfontfeatures{Scale=MatchLowercase}
  \defaultfontfeatures[\rmfamily]{Ligatures=TeX,Scale=1}
\fi
% Use upquote if available, for straight quotes in verbatim environments
\IfFileExists{upquote.sty}{\usepackage{upquote}}{}
\IfFileExists{microtype.sty}{% use microtype if available
  \usepackage[]{microtype}
  \UseMicrotypeSet[protrusion]{basicmath} % disable protrusion for tt fonts
}{}
\makeatletter
\@ifundefined{KOMAClassName}{% if non-KOMA class
  \IfFileExists{parskip.sty}{%
    \usepackage{parskip}
  }{% else
    \setlength{\parindent}{0pt}
    \setlength{\parskip}{6pt plus 2pt minus 1pt}}
}{% if KOMA class
  \KOMAoptions{parskip=half}}
\makeatother
\usepackage{xcolor}
\IfFileExists{xurl.sty}{\usepackage{xurl}}{} % add URL line breaks if available
\IfFileExists{bookmark.sty}{\usepackage{bookmark}}{
\usepackage{hyperref}
}
\hypersetup{
  pdftitle={Préparation à l'oral du Bac},
  pdfauthor={Terminale ISN Lycée du Parc},
  colorlinks=true,
  linkcolor=Maroon,
  filecolor=Maroon,
  citecolor=Blue,
  urlcolor=Blue,
  pdfcreator={LaTeX via pandoc}}
\urlstyle{same} % disable monospaced font for URLs
\usepackage[top=20mm,left=20mm,right=20mm,heightrounded]{geometry}
\usepackage{listings}
\newcommand{\passthrough}[1]{#1}
\lstset{defaultdialect=[5.3]Lua}
\lstset{defaultdialect=[x86masm]Assembler}
\usepackage{longtable,booktabs}
% Correct order of tables after \paragraph or \subparagraph
\usepackage{etoolbox}
\makeatletter
\patchcmd\longtable{\par}{\if@noskipsec\mbox{}\fi\par}{}{}
\makeatother
% Allow footnotes in longtable head/foot
\IfFileExists{footnotehyper.sty}{\usepackage{footnotehyper}}{\usepackage{footnote}}
\makesavenoteenv{longtable}
\setlength{\emergencystretch}{3em} % prevent overfull lines
\providecommand{\tightlist}{%
  \setlength{\itemsep}{0pt}\setlength{\parskip}{0pt}}
\setcounter{secnumdepth}{5}

\title{Préparation à l'oral du Bac}
\author{Terminale ISN Lycée du Parc}
\date{}

%%%jolis boites

\usepackage{fancybox, graphicx}



%%%%%%%%%%%%%%%%Packages et Macros Frederic%%%%%%%%%%%%%%%%%%%%%%%%%%%%%


%%%%Insertion de liens hypertextes %%%%

            
%%%%%%%%%%PSTricks%%%%%%%%%%%%

\usepackage{pstricks,pst-plot,pst-text,pst-tree,pst-eps,pst-fill,pst-node,pst-math,pstricks-add,pst-xkey,pst-eucl}


%%%%%%%Tikz%%%%%%%%%%%%%%%
\usepackage{pgf,tikz,tkz-tab}
% Pour les tableaux de signes ou de variations avec tkz-tab voir https://zestedesavoir.com/tutoriels/439/des-tableaux-de-variations-et-de-signes-avec-latex/#1-13389_tikz-un-package-qui-en-a-dans-le-ventre
\usetikzlibrary{arrows}
\usetikzlibrary{shapes.geometric}
\usetikzlibrary{shapes.geometric}
\usetikzlibrary{petri}
\usetikzlibrary{decorations}
\usetikzlibrary{arrows}
\usetikzlibrary{math}
 %Variables must be declared in a tikzmath environment but
       % can be used outside
%       \tikzmath{int \n; \n = 508; \x1 = 1; \y1 =1; 
%                   %computations are also possible
%                    \x2 = \x1 + 1; \y2 =\y1 +3; } 


%%%%%%%%%%%%%%%%%%%%%%%%%%%%%%%%%%%%%%%%
%%%%%%%%%%%Commandes Tikz Perso%%%%%%%%%%%%%%%

% Définition des nouvelles options xmin, xmax, ymin, ymax
% Valeurs par défaut : -3, 3, -3, 3
\tikzset{
xmin/.store in=\xmin, xmin/.default=-3, xmin=-3,
xmax/.store in=\xmax, xmax/.default=3, xmax=3,
ymin/.store in=\ymin, ymin/.default=-3, ymin=-3,
ymax/.store in=\ymax, ymax/.default=3, ymax=3,
}
% Commande qui trace la grille entre (xmin,ymin) et (xmax,ymax)
\newcommand {\grille}[2]
{\draw[help lines,black, thick] (\xmin,\ymin) grid[xstep=#1, ystep=#2] (\xmax,\ymax);}
% Commande \axes
\newcommand {\axes} {
\draw[->,very thick] (\xmin,0) -- (\xmax,0);
\draw[->,very thick] (0,\ymin) -- (0,\ymax);
\draw (0.95*\xmax, 0) node[above] {};
\draw (0, 0.95*\ymax) node[left] {};
}
% Commande qui limite l?affichage à (xmin,ymin) et (xmax,ymax)
\newcommand {\fenetre}
{\clip (\xmin,\ymin) rectangle (\xmax,\ymax);}

%Exemple d'utilisation

%\begin{center}
%\begin{tikzpicture} [xmin=-2,xmax=2,ymin=0,ymax=5]
%\grille{1} \axes \fenetre
%\draw plot[smooth] (\x,\x^2);
%\end{tikzpicture}
%\end{center}

%style pour la perspective cavalière française
%voir Tikz pour l'impatient page 68
\tikzset{math3d/.style=
{x= {(-0.353cm,-0.353cm)}, z={(0cm,1cm)},y={(1cm,0cm)}}}

%%%%%%%Symbole pour code calculatrice%%%%%%

%Flèche remplie pour défilement de menu

\newcommand{\flechefillright}{
\begin{tikzpicture}[scale=0.15] \fill (0,0)--(2,1)--(0,2)--cycle;
\end{tikzpicture}}

%%%%%%%%%%%%%Symboles pour calculatrice Casio%%%%
\newcommand{\execasio}{\Pisymbol{psy}{191}} %Retour chariot
\newcommand{\dispcasio}{\begin{pspicture}(.1,.1)\pspolygon*(.1,0)(.1,.1)\end{pspicture}} %Triangle « Disp »
\newcommand{\dispcasiotikz}{\begin{tikzpicture}[scale=0.2]
\fill (0,0) -- (1,0) -- (1,1) -- cycle;
\end{tikzpicture}} %Triangle « Disp »
%

%Fleche entre deux lignes, d'apres 'un bon petit' : http://forum.mathematex.net/latex-f6/fleches-entre-deux-lignes-pour-resolution-d-equation-t10283.html#p99817
\newcommand\addnode[1]{\Rnode{#1}{}}
\newcommand\linknode[3]{\ncbar[angleA=0,angleB=0,nodesep=1ex,arm=10ex,offset=-2pt]{->}{#1}{#2}\Aput{\vphantom{x}#3}}


%%Commande pour touche de calculatrice

\newcommand\tc[1]{%
{
\begin{tikzpicture}
\node[draw,rectangle,rounded corners=3pt] (P) at (0,0){#1};
\end{tikzpicture}
}
}

%%%%%%%%%%%%%%%%%%%%%%%%%%%%%%%%%%%%%%%%
%%%%%%%%%%%Fin Commandes Tikz%%%%%%%%%%%%%%%


%%%%%%%%%%%%Specifiques%%%%%%%%%%%
\usepackage{wrapfig}
%pour insérer une figure à droite ou à gauche d'un texte
%\begin{wrapfigure}[nb lignes]{placement l,r,c,i(inside),o(outside)}[overhang]{width}
%ce package fonctionne mal à proximité des listes
%%%%%%%%%%%%%%%%%%%%%%%%%%%%%%%%%%%%%

%%%%%Environnements et symboles spéciaux pour faire joli%%%%%%

%%%Bclogo, pour des environnements + jolis avec insertion de logo%%%%
%Dépendances de  bclogo
\usepackage{xkeyval}  
\usepackage{etoolbox}
\usepackage{ifpdf}
\usepackage[framemethod=tikz]{mdframed}
\usepackage[tikz]{bclogo}

%\newcommand\bcpython{\includegraphics[width=17pt]{/home/fjunier/Maths/python-logo.eps}}
\newcommand\bcpython{\includegraphics[width=17pt]{/home/fjunier/Maths/python-logo.png}}
%\newcommand\bcpython{\includegraphics[width=17pt]{/home/frederic/Maths/python-logo.png}}

%% Framed
\usepackage{framed}  %Le package « framed» Crée 3 nouveaux environnements, qui se comportent comme des minipage de largeur \linewidth, mais permettant en plus de se casser entre plusieurs pages.     * framed : avec un cadre autour;     * shaded : avec un fonc coloré (il faut définir la couleur shadecolor);     * leftbar : avec une barre le long du côté gauche.

%%%%%%%%%%%%%%%%%%%Présentation de codes sources%%%%%%%%%%%%%%%%%
\usepackage{listings}
%On utilise l?environnement lstlisting pour insérer
%un code source.
%En plus de l?environnement lstlisting, on peut également utiliser la
%commande \lstinline qui fonctionne comme la commande \verb, en ce
%sens qu?on peut utiliser n?importe quel caractère comme délimiteur. Enfin,
%la commande \lstinputlisting permet de charger un code source depuis
%un fichier externe.
%Il y a deux manières de préciser des options : soit via l?option de l?envi-
%ronnement ou de la commande, soit en utilisant la commande \lstset
%qui permet de définir des options de manière globale.

\lstset{ %
  language=Python,                % the language of the code
  basicstyle=\ttfamily,           % the size of the fonts that are used for the code
  %numbers=left,                   % where to put the line-numbers
  numberstyle=\tiny,  % the style that is used for the line-numbers
  %stepnumber=2,                   % the step between two line-numbers. If it's 1, each line 
                                  % will be numbered
  %numbersep=5pt,                  % how far the line-numbers are from the code
  backgroundcolor=\color{white},      % choose the background color. You must add \usepackage{color}
  showspaces=false,               % show spaces adding particular underscores
  showstringspaces=false,         % underline spaces within strings
  showtabs=false,                 % show tabs within strings adding particular underscores
  frame=single,                   % adds a frame around the code
  rulecolor=\color{black},        % if not set, the frame-color may be changed on line-breaks within not-black text (e.g. comments (green here))
  tabsize=4,                      % sets default tabsize to 2 spaces
  captionpos=b,                   % sets the caption-position to bottom
  breaklines=true,                % sets automatic line breaking
  breakatwhitespace=false,        % sets if automatic breaks should only happen at whitespace
  %title=\lstname,                   % show the filename of files included with \lstinputlisting;
                                  % also try caption instead of title
  breakindent=1cm,
  keywordstyle=\color{blue},          % keyword style
  commentstyle=\color{red},       % comment style
  %stringstyle=\ttfamily\color{green},         % string literal style
  escapeinside={\%*}{*)},            % if you want to add LaTeX within your code
  morekeywords={*,...},              % if you want to add more keywords to the set
  deletekeywords={...}              % if you want to delete keywords from the given language
  upquote=true,columns=flexible,
xleftmargin=1cm,xrightmargin=1cm,
 inputencoding=utf8,			%Les lignes qui suivent sont pour le codage utf8
  extendedchars=true,
  literate=%
            {é}{{\'{e}}}1
            {è}{{\`{e}}}1
            {ê}{{\^{e}}}1
            {ë}{{\¨{e}}}1
            {û}{{\^{u}}}1
            {ù}{{\`{u}}}1
            {â}{{\^{a}}}1
            {à}{{\`{a} }}1
            {î}{{\^{i}}}1
            {ô}{{\^{o}}}1
            {ç}{{\c{c}}}1
            {Ç}{{\c{C}}}1
            {É}{{\'{E}}}1
            {Ê}{{\^{E}}}1
            {À}{{\`{A}}}1
            {Â}{{\^{A}}}1
            {Î}{{\^{I}}}1
}

\lstdefinestyle{rond}{
  numbers=none,
  backgroundcolor=\color{gristclair},
  frameround =tttt
}

\lstdefinestyle{compil}{
  numbers=none,
  backgroundcolor=\color{gristclair}
}
%\lstset{language=Python,basicstyle=\small , frame=single,tabsize=4,showspaces=false,showtabs=false,showstringspaces=false,numbers=left,numberstyle=\tiny , extendedchars=true}



%%%%%%%%%%%%%%%%%%%%%%%%%%%%%%%%%%%%%%%%%%%%%%%%%%%%%%%%%%%%%%%%%%%%%%%%
%%%%%%%%%%%%%%%%%%%%Environnements persos%%%%%%%%%%%%%%%%%%%%%%%%%%%%%%%%
%Syntaxe :
%\newenvironment{nom}[nombre d'args][defaut]{definitions initiales}{definitions finales}
%definitions intiales sont les commandes appelées par \begin{nom}
%Definitions finales sont les commandes appelées par \end{nom}

%%%%%%%%%%%%%%%%Définitions des environnemts persostheoreme, exemple ..%%%%
%%%% Exercice avec encadré %%%%
\newcounter{exo}
\newenvironment{exercice}[1]
{\par \medskip   \addtocounter{exo}{1} \noindent  
\begin{bclogo}[arrondi =0.1,   noborder = true, logo=\bccrayon, marge=4]{~\textbf{Exercice} \textbf{\theexo} {\itshape #1} }  \par}
{
\end{bclogo}
 \par \bigskip }

%%Axiomes, Theoremes, Propriété, Définition, Methode, Preuve


\newenvironment{axiome}[1]
{\par \medskip   \begin{leftbar} \noindent \underline{\textbf{Axiome}}\hspace{0.5cm}{\itshape #1}   \vspace*{10pt} \par }
{\end{leftbar}  \par \medskip }


\newcounter{thme}
\newenvironment{theoreme}[1]
{\par \medskip  \addtocounter{thme}{1} \noindent  
\begin{bclogo}[arrondi =0.1,  ombre = true, barre=none, logo=\bcbook, marge=4]{~\textbf{Théorème} \textbf{\thethme} {\itshape #1} }   \par}
{
\end{bclogo}
 \par \bigskip}

 \newenvironment{theoremedef}[1]
{\par \medskip   \addtocounter{thme}{1} \noindent  
\begin{bclogo}[arrondi =0.1,  ombre = true, barre=none, logo=\bcbook, marge=4]{~\textbf{Théorème-Définition} \textbf{\thethme} {\itshape #1} }   \par}
{
\end{bclogo}
 \par \bigskip }
 
\newcounter{prop}
\newenvironment{propriete}[1]
{\par \medskip   \addtocounter{prop}{1} \noindent  
\begin{bclogo}[arrondi =0.1,  ombre = true, barre=none, logo=\bcbook, marge=4]{~\textbf{Propriété} \textbf{\theprop} {\itshape #1} }   \par}
{
\end{bclogo}
 \par \bigskip }


\newenvironment{corollaire}[1]
{\par \medskip   \noindent  
\begin{bclogo}[arrondi =0.1,  ombre = true, barre=none, logo=\bcbook, marge=4]{~\textbf{Corollaire} {\itshape #1} } \par }
{
\end{bclogo}
 \par \bigskip }

\newenvironment{demo}[1]
{\par \medskip   \noindent  
\begin{bclogo}[arrondi =0.1,  ombre = true, barre=zigzag, noborder = true, logo=\bcloupe, marge=0]{~\textbf{Démonstration} {\itshape #1} } \par \vspace{10pt}}
{
\end{bclogo}
 \par \bigskip }

\newcounter{activite}
\newenvironment{activite}[1]
{\par \medskip   \noindent   \addtocounter{activite}{1}
\begin{bclogo}[arrondi =0.1,   noborder = true, logo=\bcvelo, marge=4]{~\textbf{Activité} \textbf{\theactivite} {\itshape #1} }  \par}
{
\end{bclogo}
 \par \bigskip }


\newcounter{rque}
\newenvironment{remarque}
{\par \medskip    \addtocounter{rque}{1} \noindent  
\begin{bclogo}[arrondi =0.1,  ombre = true, barre=snake, noborder = true, logo=\bcinfo, marge=0]{~\textbf{Remarque} \textbf{\therque}}  \par }
{
\end{bclogo}
 \par \bigskip }

\newcounter{def}
\newenvironment{definition}[1]
{\par \medskip   \addtocounter{def}{1} \noindent  
\begin{bclogo}[arrondi =0.1,  ombre = true, barre=none, logo=\bcbook, marge=4]{~\textbf{Définition} \textbf{\thedef} {\itshape #1} }  \par}
{
\end{bclogo}
 \par \bigskip }

\newcounter{exple}
\newenvironment{exemple}[1]
{\par \medskip   \addtocounter{exple}{1} \noindent  
\begin{bclogo}[arrondi =0.1,   noborder = true, logo=\bccrayon, marge=4]{~\textbf{Exemple} \textbf{\theexple} {\itshape #1} }  \par}
{
\end{bclogo}
 \par \bigskip }




\newcounter{alg}
\newenvironment{algorithme}[1]
{\par \medskip   \addtocounter{alg}{1} \noindent  
 \begin {bclogo}[noborder = true, barre=zigzag,logo=\bcpython, marge=4] {~\textbf{Algorithmique} \textbf{\thealg} {\itshape #1} }  \par}
{
\end{bclogo}
 \par \bigskip }

\newcounter{prog}
\newenvironment{programme}[1]
{\par \medskip   \addtocounter{prog}{1} \noindent  
 \begin {bclogo}[noborder = true, barre=zigzag,logo=\bcpython, marge=4] {~\textbf{Programme} \textbf{\theprog} {\itshape #1} }  \par  \bigskip}
{
\end{bclogo}
 \par \bigskip }
 
\newcounter{logi}
\newenvironment{logique}[1]
{\par \medskip   \addtocounter{logi}{1} \noindent  
 \begin {bclogo}[noborder = true, barre=zigzag,logo=\bclampe, marge=4] {~\textbf{Logique} \textbf{\thelogi} {\itshape #1} }  \par}
{
\end{bclogo}
 \par \bigskip }


\newenvironment{methode}[1]
{\par \medskip    \noindent  
 \begin {bclogo}[arrondi =0.1,logo=\bcoutil, marge=4,noborder = true] {~\textbf{Méthode}   {\itshape #1} }  \par}
{
\end{bclogo}
 \par \bigskip }


\newcounter{histo}
\newenvironment{histoire}[1]
{\par \medskip   \addtocounter{histo}{1} \noindent  
 \begin {bclogo}[couleur = blue!10 , arrondi =0.1,logo=\bchorloge, marge=4] {~\textbf{Histoire} \textbf{\thehisto} {\itshape #1} }  \par}
{
\end{bclogo}
 \par \bigskip }




%Environnement contenu pour un document présentant une progression annuelle
\newenvironment{contenu}
{\par \medskip   \begin {bclogo}[ noborder = true,logo=\bccrayon] \noindent {\large \textbf{Contenu de la séance}} \vspace*{10pt} \par  }
{\end{bclogo}  \par \medskip }

%Environnement programme pour un document présentant une progression annuelle
%\newenvironment{programme}
%{\par \medskip   \begin {bclogo}[ noborder = true, barre=zigzag,logo=\bcinfo] \noindent {\large \textbf{Programme officiel}} \vspace*{10pt} \par  }
%{\end{bclogo}  \par \medskip }

%Environnement programme pour un document présentant une progression annuelle
\newenvironment{ressource}
{\par \medskip   \begin {bclogo}[ noborder = true,logo=\bcbook] \noindent {\large \textbf{Ressources}}\\vspace*{10pt} \par }
{\end{bclogo}  \par \medskip }




%%%%%%%%%%%%%%%%%%Maths divers%%%%%%%%%%%%%%%%%%%%%%%%%
%%%%%%%%%%%%%Nombres%%%%%%%%%%%%%%%%

%Ensemble prive de...
%\newcommand{\prive}{\boi}%{\backslash}

%Ensembles de nombres%%%%%%%%%%%%%%%%%
\newcommand{\R}{\mathbb{R}}
\newcommand{\N}{\mathbb{N}}
\newcommand{\D}{\mathbb{D}}
\newcommand{\Z}{\mathbb{Z}}
\newcommand{\Q}{\mathbb{Q}}
%\newcommand{\C}{\mathbb{C}}
\newcommand{\df}{~\ensuremath{]0;+\infty[}~}
\newcommand{\K}{\mathbb{K}}

%%%%%%%%Arithmetique%%%%%%%%%%
%PGCD, PPCM
\newcommand{\PGCD}{\mathop{\rm PGCD}\nolimits}
\newcommand{\PPCM}{\mathop{\rm PPCM}\nolimits}

%Intervalles
\newcommand{\interoo}[2]{]#1\, ;\, #2[}
\newcommand{\Interoo}[2]{\left]#1\, ;\, #2\right[}
\newcommand{\interof}[2]{]#1\, ;\, #2]}
\newcommand{\Interof}[2]{\left]#1\, ;\, #2\right]}
\newcommand{\interfo}[2]{[#1\, ;\, #2[}
\newcommand{\Interfo}[2]{\left[#1\, ;\, #2\right[}
\newcommand{\interff}[2]{[#1\, ;\, #2]}
\newcommand{\Interff}[2]{\left[#1\, ;\, #2\right]}
%\newcommand\interentiers #1#2{[\! [#1\, ;\, #2]\! ]}
\newcommand{\interentiers}[2]{\llbracket #1\, ;\, #2\rrbracket}
%


%%%%%%%%%%%%%%Nombres complexes%%%%%

\newcommand{\ic}{\text{i}}
%\newcommand{\I}{\text{i}}
\newcommand{\im}[1]{\text{Im}\left(#1\right)}
\newcommand{\re}[1]{\text{Re}\left(#1\right)}
\newcommand{\Arg}[1]{\text{arg}\left(#1\right)}
\newcommand{\Mod}[1]{\left[#1\right]}
%Parties entière, réelle, imaginaire, nombre i
\newcommand{\ent}[1]{\text{E}\left(#1\right)}
\renewcommand{\Re}{\mathop{\rm Re}\nolimits}
\renewcommand{\Im}{\mathop{\rm Im}\nolimits}
\renewcommand{\i}{\textrm{i}}

%%%%%%%%%%%Probabilites et statistiques%%%%%
\newcommand{\loibinom}[2]{\mathcal{B}\left(#1\ ; \ #2 \right)}
\newcommand{\loinorm}[2]{\mathcal{N}\left(#1\ ; \ #2 \right)}
\newcommand{\loiexp}[1]{\mathcal{E}\left(#1\right)}
\newcommand{\proba}[1]{\mathbb{P}\big(#1\big)}
\newcommand{\probacond}[2]{\mathbb{P}_{#2}\big(#1\big)}
\newcommand{\esperance}[1]{\mathbb{E}\left(#1\right)}
\newcommand{\variance}[1]{\mathbb{V}\left(#1\right)}
\newcommand{\ecart}[1]{\sigma\left(#1\right)}
\newcommand{\dnormx}{\frac{1}{\sqrt{2\pi}} \text{e}^{-\frac{x^2}{2}}}
\newcommand{\dnormt}{\frac{1}{\sqrt{2\pi}} \text{e}^{-\frac{t^2}{2}}}

%Covariance
\newcommand{\cov}{\mathop{\rm cov}\nolimits}
%


%%%%%%%%%%Analyse%%%%%%%%%%%

%%%%%%%%%%%Courbe%%%%%%%%%%%%
\newcommand{\courbe}[1]{\ensuremath{\mathcal{C}_{#1}}}

%%%%%%%Fonction exponentielle%%%%%
\newcommand{\fe}{~fonction exponentielle~}
\newcommand{\e}{\text{e}}

%Fonction cotangente
\newcommand{\cotan}{\mathop{\rm cotan}\nolimits}
%%%%%%%%%%%%%%%%%%%%%%%%%%%%%%%%%%%%%%%%%
%
%Fonctions hyperboliques
\newcommand{\ch}{\mathop{\rm ch}\nolimits}
\newcommand{\sh}{\mathop{\rm sh}\nolimits}


%%%%%%%%%%%%%%Limites%%%%%%
\newcommand{\limite}[2]{\lim\limits
_{x \to #1} #2}
\newcommand{\limitesuite}[1]{\lim\limits
_{n \to +\infty} #1}
\newcommand{\limiteg}[2]{\lim\limits
_{\substack{x \to #1 \\ x < #1 }} #2}
\newcommand{\limited}[2]{\lim\limits
_{\substack{x \to #1 \\ x > #1 }} #2}

%%%%%%%%%%Continuité%%%%%%%%%%%
\newcommand{\TVI}{théorème des valeurs intermédiaires}

%%%%%%%%%%%Suites%%%%%%%%%%%%
\newcommand{\suite}[1]{\ensuremath{\left(#1_{n}\right)}}
\newcommand{\Suite}[2]{\ensuremath{\left(#1\right)_{#2}}}
%

%%%%%%%%%%%%%%%Calcul intégral%%%%%%
\newcommand{\dx}{\ensuremath{\text{d}x}}		% dx
\newcommand{\dt}{\ensuremath{\text{d}t}}		% dt
\newcommand{\dtheta}{\ensuremath{\text{d}\theta}}		% dtheta
\newcommand{\dy}{\ensuremath{\text{d}y}}		% dy
\newcommand{\dq}{\ensuremath{\text{d}q}}		% dq

%%%Intégrale%%%
\newcommand{\integralex}[3]{\int_{#1}^{#2} #3 \ \dx}
\newcommand{\integralet}[3]{\int_{#1}^{#2} #3 \ \dt}
\newcommand{\integraletheta}[3]{\int_{#1}^{#2} #3 \ \dtheta}

%%%%%Equivalent%%
\newcommand{\equivalent}[1]{\build\sim_{#1}^{}}

%o et O%%%%
\renewcommand{\o}[2]{\build o_{#1\to #2}^{}}
\renewcommand{\O}[2]{\build O_{#1\to #2}^{}}



%%%%%%%%%%%%%%%Geometrie%%%%%%%%%%%%%%%%%%%%%%%

%%%%%%%%%%%%%%%Reperes%%%%%%%%%%%%%%
\def\Oij{\ensuremath{\left(\text{O},~\vect{\imath},~\vect{\jmath}\right)}}
\def\Oijk{\ensuremath{\left(\text{O},~\vect{\imath},~ \vect{\jmath},~ \vect{k}\right)}}
\def\Ouv{\ensuremath{\left(\text{O},~\vect{u},~\vect{v}\right)}}
\renewcommand{\ij}{(\vec\imath\, ;\vec\jmath\,)}
\newcommand{\ijk}{(\vec\imath\, ;\vec\jmath\, ;\vec k\,)}
\newcommand{\OIJ}{(O\,;\, I\,;\, J\,)}
\newcommand{\repere}[3]{\big(#1\, ;\,\vect{#2} ;\vect{#3}\big)}
\newcommand{\reperesp}[4]{\big(#1\, ;\,\vect{#2} ;\vect{#3} ;\vect{#4}\big)}

%%%%%%%%%Coordonnees%%%%%%%%%%%%%%
\newcommand{\coord}[2]{(#1\, ;\, #2)}
\newcommand{\bigcoord}[2]{\big(#1\, ;\, #2\big)}
\newcommand{\Coord}[2]{\left(#1\, ;\, #2\right)}
\newcommand{\coordesp}[3]{(#1\, ;\, #2\, ;\, #3)}
\newcommand{\bigcoordesp}[3]{\big(#1\, ;\, #2\, ;\, #3\big)}
\newcommand{\Coordesp}[3]{\left(#1\, ;\, #2\, ;\, #3\right)}
\newcommand{\Vcoord}[3]{\begin{pmatrix} #1 \\ #2 \\ #3 \end{pmatrix}}
%Symboles entre droites
%\newcommand{\paral}{\sslash}
\newcommand{\paral}{\mathop{/\!\! /}}
%

%%%%%%%%%Produit scalaire, Angles%%%%%%%%%%
\newcommand{\scal}[2]{\vect{#1} \, \cdot \, \vect{#2}}
\newcommand{\Angle}[2]{\left(\vect{#1} \, , \, \vect{#2}\right)}
\newcommand{\Anglegeo}[2]{\left(\widehat{\vect{#1} \, ; \, \vect{#2}}\right)}
\renewcommand{\angle}[1]{\widehat{#1}}
\newcommand{\anglevec}[2]{\left(\vect {#1}\, ,\,\vect {#2} \right)}
\newcommand{\anglevecteur}[2]{(#1\, , \, #2)}
\newcommand{\Anglevec}[2]{(\vecteur{#1}\, ,\,\vecteur{#2})}
\newcommand{\prodscal}[2]{#1 \, \cdot \, #2}
%


%Arc
%\newcommand{\arc}[1]{\wideparen{#1}}
\newcommand{\arcoriente}[1]{\overset{\curvearrowright}{#1}}
%
%


%%%%%%%%%%%%%%%Normes%%%%%%%%%%%%%%%%
\newcommand{\norme}[1]{\left\| #1\right\|}
\newcommand{\normebis}[1]{\delim{2pt}{\|}{9pt}\! #1\delim{2pt}{\|}{9pt}}
\newcommand{\normetriple}[1]{\left |\kern -.07em\left\| #1\right |\kern -.07em\right\|}
\newcommand{\valabs}[1]{\big| \, #1 \, \big|}
%

%%%%%%%%%%%%%%%%%%%%%%%%%%%Degré%%%%%%
%\newcommand{\Degre}{\ensuremath{^\circ}}
%La commande \degre est déjà définie dans le package babel

%%%%%%%%%%Vecteurs%%%%%%%%%%%
\newcommand{\vect}[1]{\mathchoice%
{\overrightarrow{\displaystyle\mathstrut#1\,\,}}%
{\overrightarrow{\textstyle\mathstrut#1\,\,}}%
{\overrightarrow{\scriptstyle\mathstrut#1\,\,}}%
{\overrightarrow{\scriptscriptstyle\mathstrut#1\,\,}}}



%%%%%%%%%%%%%Algebre%%%%%%%%%%%%%%%


%%%%%%%%%%Systemes%%%%%%%%%%%
%Systemes
\newcommand{\sys}[2]{
\left\lbrace
 \begin{array}{l}
  \negthickspace\negthickspace #1\\
  \negthickspace\negthickspace #2\\
 \end{array}
\right.\negthickspace\negthickspace}
\newcommand{\Sys}[3]{
\left\lbrace
 \begin{array}{l}
  #1\\
  #2\\
  #3\\
 \end{array}
\right.}
\newcommand{\Sysq}[4]{
\left\lbrace
 \begin{array}{l}
  #1\\
  #2\\
  #3\\
  #4\\
 \end{array}
\right.}
%
%

%%%%%%%%%%%%%%%%Matrices%%%%%%%%%%%%%%%%%%
%Comatrice
\newcommand{\com}{\mathop{\rm com}\nolimits}
%
%
%Trace
\newcommand{\tr}{\mathop{\rm tr}\nolimits}
%
%
%Transposee
\newcommand{\transposee}[1]{{\vphantom{#1}}^t\negmedspace #1}
%
%
%Noyau
\newcommand{\Ker}{\mathop{\rm Ker}\nolimits}
%
%

%
%Matrices
\newcommand{\Mn}{\mathcal M_n}
\newcommand{\matrice}[4]{
\left(
 \begin{array}{cc}
  #1 & #2 \\
  #3 & #4
 \end{array}
\right)}

\newcommand{\Matrice}[9]{
\left(
 \begin{array}{ccc}
  #1 & #2 & #3\\
  #4 & #5 & #6\\
  #7 & #8 & #9
 \end{array}
\right)}
\newcommand{\Vect}[3]{
\left(\negmedspace
 \begin{array}{c}
  #1\\
  #2\\
  #3
 \end{array}\negmedspace
\right)}
\newcommand{\Ideux}{\matrice{1}{0}{0}{1}}
\newcommand{\Itrois}{\Matrice{1}{0}{0}{0}{1}{0}{0}{0}{1}}
%
%
%Determinants
\newcommand{\determinant}[4]{
\left|
 \begin{array}{cc}
  #1 & #2 \\
  #3 & #4
 \end{array}
\right|}
\newcommand{\Determinant}[9]{
\left|
 \begin{array}{ccc}
  #1 & #2 & #3\\
  #4 & #5 & #6\\
  #7 & #8 & #9
 \end{array}
\right|}

\begin{document}
\maketitle

\renewcommand*\contentsname{Table des matières}
{
\hypersetup{linkcolor=}
\setcounter{tocdepth}{3}
\tableofcontents
}
\hypertarget{cruxe9dits}{%
\section{Crédits}\label{cruxe9dits}}

Plusieurs exercices sont tirés des questions proposées sur le site
\url{https://genumsi.inria.fr/}.

\hypertarget{programmation-en-python}{%
\section{Programmation en Python}\label{programmation-en-python}}

\begin{exercice}{}

\emph{Auteur : Aude Durand Senter, question n°178 Genumsi}

On a saisi le code suivant :

\begin{lstlisting}[language=Python]
def mystere(nombre) :
    while nombre > 5 :
        nombre = nombre - 5
return nombre
\end{lstlisting}

Quelle affirmation est vraie dans celles proposées ci-dessous ?

Réponses :

\begin{enumerate}
\def\labelenumi{\arabic{enumi}.}
\tightlist
\item
  On sort de la boucle while dès que
  \passthrough{\lstinline!nombre > 5!}
\item
  On sort de la boucle while dès que
  \passthrough{\lstinline!nombre < 5!}
\item
  On sort de la boucle while dès que
  \passthrough{\lstinline!nombre <= 5!}
\item
  On continue la boucle tant que \passthrough{\lstinline!nombre <= 5!}
\end{enumerate}

\end{exercice}

\begin{exercice}{}

\emph{Auteur : Nicolas Revéret}

On considère le programme ci-dessous :

\begin{lstlisting}[language=Python]
a = 8
b = 5
a = a + b
b = a - b
a = a - b
\end{lstlisting}

Quelles sont les valeurs des variables \passthrough{\lstinline!a!} et
\passthrough{\lstinline!b!} à la fin du programme ?

\end{exercice}

\begin{exercice}{}

\emph{Auteur : Nicolas Revéret}

On souhaite que l'utilisateur saisisse une valeur entière comprise entre
1 et 20 (inclus).

Quel code est correct ?

Réponses :

\begin{enumerate}
\def\labelenumi{\arabic{enumi}.}
\tightlist
\item
  Réponse 1
\end{enumerate}

\begin{lstlisting}[language=Python]
a = int ( input("Saisir un nombre entre 1 et 20") )
while (a < 1 and a > 20) :
      a = int ( input("Saisir un nombre entre 1 et 20") )
\end{lstlisting}

\begin{enumerate}
\def\labelenumi{\arabic{enumi}.}
\setcounter{enumi}{1}
\tightlist
\item
  Réponse 2
\end{enumerate}

\begin{lstlisting}[language=Python]
while (a < 1 and a > 20) :
    a = int ( input("Saisir un nombre entre 1 et 20") )
\end{lstlisting}

\begin{enumerate}
\def\labelenumi{\arabic{enumi}.}
\setcounter{enumi}{2}
\tightlist
\item
  Réponse 3
\end{enumerate}

\begin{lstlisting}[language=Python]
a = int ( input("Saisir un nombre entre 1 et 20") )
while (a < 1 or a > 20) :
    a = int ( input("Saisir un nombre entre 1 et 20") )
\end{lstlisting}

\begin{enumerate}
\def\labelenumi{\arabic{enumi}.}
\setcounter{enumi}{3}
\tightlist
\item
  Réponse 4
\end{enumerate}

\begin{lstlisting}[language=Python]
a = int ( input("Saisir un nombre entre 1 et 20") )
while (a >= 1 and a >= 20) :
    a = int ( input("Saisir un nombre entre 1 et 20") )
\end{lstlisting}

\end{exercice}

\begin{exercice}{}

\emph{Auteur : Nicolas Revéret}

On a saisi le code suivant :

\begin{lstlisting}[language=Python]
a = 12
for i in range(3) :
  a = a * 2
  a = a - 10
\end{lstlisting}

Quelle est la valeur de la variable \passthrough{\lstinline!a!} après
l'exécution du code ?

\end{exercice}

\begin{exercice}{}

Pour les questions suivantes, il faut utiliser obligatoirement deux
boucles imbriquées et au maximum deux fois la fonction
\passthrough{\lstinline!print!}.

\begin{enumerate}
\def\labelenumi{\arabic{enumi}.}
\tightlist
\item
  Ecrire un script Python qui produit l'affichage 1 ci-dessous.
\item
  Ecrire un script Python qui produit l'affichage 2 ci-dessous.
\end{enumerate}

\textbf{Affichage 1}

\begin{lstlisting}[language=Python]
0
01
012
0123
01234
\end{lstlisting}

\textbf{Affichage 2}

\begin{lstlisting}[language=Python]
01234
12345
23456
34567
45678
\end{lstlisting}

\end{exercice}

\begin{exercice}{}

Pour le diplôme du baccalauréat, si on note \(m\) la moyenne du
candidat, quatre mentions sont possibles : \emph{Passable} si
\(10 \leqslant m < 12\), \emph{Assez bien} si \(12 \leqslant m < 14\),
\emph{Bien} si \(14 \leqslant m < 16\) et \emph{Très bien} sinon.
Recopier et compléter le script Python ci-dessous pour qu'il affiche la
mention d'un candidat admis (on suppose sa moyenne supérieure ou égale à
10).

\begin{lstlisting}[language=Python]
m = float(input('Moyenne du candidat ? '))
if 10 <= m < 12:
    print("Passable")
#to be completed
\end{lstlisting}

\end{exercice}

\begin{exercice}{}

\emph{Auteur : Nicolas Revéret}

On considère le code suivant :

\begin{lstlisting}[language=Python]
def f(tab):
  for i in range(len(tab)//2):
    tab[i],tab[-i-1] = tab[-i-1],tab[i]
\end{lstlisting}

Après les lignes suivantes :

\begin{lstlisting}[language=Python]
tab = [2,3,4,5,7,8]
f(tab)
\end{lstlisting}

Quelle est la valeur de la variable \passthrough{\lstinline!tab!} ?

Réponses :

\begin{enumerate}
\def\labelenumi{\arabic{enumi}.}
\tightlist
\item
  \passthrough{\lstinline![2,3,4,5,7,8]!}
\item
  \passthrough{\lstinline![5,7,8,2,3,4]!}
\item
  \passthrough{\lstinline![8,7,5,4,3,2]!}
\item
  \passthrough{\lstinline![4,3,2,8,7,5]!}
\end{enumerate}

\end{exercice}

\begin{exercice}{}

\emph{Auteur : Christophe Beasse}

Que contient la variable \passthrough{\lstinline!a!} si on exécute ce
script ?

\begin{lstlisting}[language=Python]
def carre(val):
  return val*val

def inc(val):
  return val + 1

a = carre(inc(3.0))
\end{lstlisting}

\begin{enumerate}
\def\labelenumi{\arabic{enumi}.}
\tightlist
\item
  9.0
\item
  12.0
\item
  10.0
\item
  16.0
\end{enumerate}

\end{exercice}

\begin{exercice}{}

\emph{Auteur : Christophe Beasse}

Soit la liste suivante :
\passthrough{\lstinline!liste\_pays = ["France","Allemagne","Italie","Belgique"]!}

Que renvoie : \passthrough{\lstinline!liste\_pays[2]!} ?

Réponses :

\begin{enumerate}
\def\labelenumi{\arabic{enumi}.}
\tightlist
\item
  \passthrough{\lstinline!"France"!}
\item
  \passthrough{\lstinline!"Allemagne"!}
\item
  \passthrough{\lstinline!"Italie"!}
\item
  \passthrough{\lstinline!"Belgique"!}
\end{enumerate}

\end{exercice}

\begin{exercice}{}

\emph{Auteur : Christophe Beasse}

Quelle est le résultat de :
\passthrough{\lstinline![ (a,b) for a in range(3) for b in range(a) ]!}
?

Réponses :

\begin{enumerate}
\def\labelenumi{\arabic{enumi}.}
\tightlist
\item
  \passthrough{\lstinline![(1,0),(2,1),(2,1)]!}
\item
  \passthrough{\lstinline![(1,0),(2,1),(3,2)]!}
\item
  \passthrough{\lstinline![(1,0),(2,0),(2,1)]!}
\item
  \passthrough{\lstinline![(0,0),(1,1),(2,2)]!}
\end{enumerate}

\end{exercice}

\begin{exercice}{}

Ecrire une fonction \passthrough{\lstinline!moyenne(t)!} qui prend en
argument un tableau de nombres t et qui retourne sa moyenne
arithmétique.

\end{exercice}

\begin{exercice}{}

\emph{Auteur : Nicolas Revéret}

On dispose d'un tableau d'entiers, ordonné en ordre croissant.

On désire connaître le nombre de valeurs distinctes contenues dans ce
tableau.

Quelle est la fonction qui ne convient pas ?

Réponses :

\begin{enumerate}
\def\labelenumi{\arabic{enumi}.}
\tightlist
\item
  Réponse 1
\end{enumerate}

\begin{lstlisting}[language=Python]
def compte(t):
    cpt = 1
    for i in range(1,len(t)):
        if t[i] != t[i-1]:
            cpt = cpt + 1
    return cpt
\end{lstlisting}

\begin{enumerate}
\def\labelenumi{\arabic{enumi}.}
\setcounter{enumi}{1}
\tightlist
\item
  Réponse 2
\end{enumerate}

\begin{lstlisting}[language=Python]
    def compte(t):
      cpt = 0
      for i in range(0,len(t)-1):
        cpt = cpt + int(t[i] != t[i+1])
      return cpt
\end{lstlisting}

\begin{enumerate}
\def\labelenumi{\arabic{enumi}.}
\setcounter{enumi}{2}
\tightlist
\item
  Réponse 3
\end{enumerate}

\begin{lstlisting}[language=Python]
    def compte(t):
      cpt = 0
      for i in range(0,len(t)-1):
        if t[i] != t[i+1]:
          cpt = cpt + 1
      return cpt+1
\end{lstlisting}

\end{exercice}

\begin{exercice}{}

\emph{Auteur : Eric Rougier}

Quel est le résultat de l'évaluation de l'expression Python suivante ?

\passthrough{\lstinline![2 ** n for n in range(4)]!}

Réponses :

\begin{enumerate}
\def\labelenumi{\arabic{enumi}.}
\tightlist
\item
  \passthrough{\lstinline![0, 2, 4, 6, 8]!}
\item
  \passthrough{\lstinline![1, 2, 4, 8]!}
\item
  \passthrough{\lstinline![0, 1, 4, 9]!}
\item
  \passthrough{\lstinline![1, 2, 4, 8, 16]!}
\end{enumerate}

\end{exercice}

\begin{exercice}{}

\emph{Auteur : Germain Becker, question n°326 Genumsi} Quel est le
tableau \passthrough{\lstinline!t!} construit par les instructions
suivantes ?

\begin{lstlisting}[language=Python]
tab = [1, 2, -3, 7, 4, 10, -1, 0]
t = [e for e in tab if e >= 0]
\end{lstlisting}

Réponses :

\begin{enumerate}
\def\labelenumi{\arabic{enumi}.}
\tightlist
\item
  \passthrough{\lstinline!t = [1, 2, 7, 4, 10, 0]!}
\item
  \passthrough{\lstinline!t = [e, e, e, e, e, e]!}
\item
  \passthrough{\lstinline!t = [1, 2, 7, 4, 10]!}
\item
  \passthrough{\lstinline!t = [-3, -1, 0]!}
\end{enumerate}

\end{exercice}

\begin{exercice}{}

\emph{Auteur : Germain Becker, question n°339 Genumsi}

On considère le tableau t suivant.

\passthrough{\lstinline!t = [[1, 2, 3], [2, 3, 4], [3, 4, 5], [4, 5, 6]]!}

Quelle est la valeur de t{[}1{]}{[}2{]} ?

Réponses :

\begin{enumerate}
\def\labelenumi{\arabic{enumi}.}
\tightlist
\item
  1
\item
  3
\item
  4
\item
  2
\end{enumerate}

\end{exercice}

\begin{exercice}{}

\emph{Auteur : Eric Rougier, question n°150 Genumsi}

Quelle est la valeur de la variable image après exécution du script
Python suivant ?

\begin{lstlisting}[language=Python]
image = [[0, 0, 0, 0], [0, 0, 0, 0], [0, 0, 0, 0], [0, 0, 0, 0]]
for i in range(4):
    for j in range(4):
        if (i+j) == 3:
            image[i][j] = 1
\end{lstlisting}

Réponses :

\begin{enumerate}
\def\labelenumi{\arabic{enumi}.}
\item
  \passthrough{\lstinline![[0, 0, 0, 0], [0, 0, 0, 0], [0, 0, 0, 0], [1, 1, 1, 1]]!}
\item
  \passthrough{\lstinline![[0, 0, 0, 1], [0, 0, 0, 1], [0, 0, 0, 1], [0, 0, 0, 1]]!}
\item
  \passthrough{\lstinline![[0, 0, 0, 1], [0, 0, 1, 0], [0, 1, 0, 0], [1, 0, 0, 0]]!}
\item
  \passthrough{\lstinline![[0, 0, 0, 1], [0, 0, 1, 1], [0, 1, 1, 1], [1, 1, 1, 1]]!}
\end{enumerate}

\end{exercice}

\begin{exercice}{}

\emph{Auteur : Nicolas Réveret, question n°379 Genumsi}

On a saisi le code suivant :

\begin{lstlisting}[language=Python]
liste = [0, 1, 2, 3]
compteur = 0
for i in range(len(liste)-1) :
    for j in range(i,len(liste)) :
        compteur += 1
\end{lstlisting}

Que contient la variable compteur à la fin de l'exécution de ce script ?

\end{exercice}

\begin{exercice}{}

Quelle est la valeur référencée par la liste L après exécution du
programme ci-dessous ?

\begin{lstlisting}[language=Python]
L = [731, 732, 734]
L[0], L[1] = L[1], L[0]
M = L
M[1] = 732
\end{lstlisting}

\begin{enumerate}
\def\labelenumi{\arabic{enumi}.}
\tightlist
\item
  {[}732, 731, 734{]}
\item
  {[}731, 732, 734{]}
\item
  {[}732, 732, 734{]}
\end{enumerate}

\end{exercice}

\begin{exercice}{}

On considère l'extrait de code suivant :

\begin{lstlisting}[language=Python]
while (a < 20) or (b > 50):
    .......
    .......
\end{lstlisting}

Quelles conditions permettent de mettre fin à cette boucle ?

\begin{enumerate}
\def\labelenumi{\arabic{enumi}.}
\item
  \textbf{Réponse 1 :} la boucle prend fin lorsque
  \passthrough{\lstinline!a < 20 ou b > 50!}
\item
  \textbf{Réponse 2 :} la boucle prend fin lorsque
  \passthrough{\lstinline!a < 20 et b > 50!}
\item
  \textbf{Réponse 3 :} la boucle prend fin lorsque
  \passthrough{\lstinline!a >= 20 ou b <= 50!}
\item
  \textbf{Réponse 4 :} la boucle prend fin lorsque
  \passthrough{\lstinline!a >= 20 et b <= 50!}
\end{enumerate}

\end{exercice}

\begin{exercice}{}

On exécute le script suivant :

\begin{lstlisting}[language=Python]
L = [12,0,8,7,3,1,5,3,8]

a = [elt for elt in L if elt < 4]
\end{lstlisting}

Quelle est la valeur de a à la fin de son exécution ?

Réponses :

\begin{enumerate}
\def\labelenumi{\arabic{enumi}.}
\item
  \textbf{Réponse 1 :} \passthrough{\lstinline![12,0,8]!}
\item
  \textbf{Réponse 2 :} \passthrough{\lstinline![12,0,8,7]!}
\item
  \textbf{Réponse 3 :} \passthrough{\lstinline![0,3,1,3]!}
\item
  \textbf{Réponse 4 :} \passthrough{\lstinline![0,3,1]!}
\end{enumerate}

\end{exercice}

\begin{exercice}{}

On définit : \passthrough{\lstinline!L = [10,9,8,7,6,5,4,3,2,1]!}.

Quelle est la valeur de \passthrough{\lstinline!L[L[3]]!} ?

Réponses :

\begin{enumerate}
\def\labelenumi{\arabic{enumi}.}
\item
  \textbf{Réponse 1 :} 3
\item
  \textbf{Réponse 2 :} 4
\item
  \textbf{Réponse 3 :} 7
\item
  \textbf{Réponse 4 :} 8
\end{enumerate}

\end{exercice}

\begin{exercice}{}

On définit une grille G remplie de 0, sous la forme d'une liste de
listes, où toutes les sous-listes ont le même nombre d'éléments.

\begin{lstlisting}[language=Python]
G = [ [0, 0, 0, ..., 0],
[0, 0, 0, ..., 0],
[0, 0, 0, ..., 0],

......
[0, 0, 0, ..., 0]]
\end{lstlisting}

On appelle \emph{hauteur} de la grille le nombre de sous-listes
contenues dans G et \emph{largeur} de la grille le nombre d'éléments
dans chacune de ces sous-listes. Comment peut-on les obtenir ?

Réponses :

\begin{enumerate}
\def\labelenumi{\arabic{enumi}.}
\tightlist
\item
  \textbf{Réponse 1:}
\end{enumerate}

\begin{lstlisting}[language=Python]
hauteur = len(G[0])
largeur = len(G)
\end{lstlisting}

\begin{enumerate}
\def\labelenumi{\arabic{enumi}.}
\setcounter{enumi}{1}
\tightlist
\item
  \textbf{Réponse 2 :}
\end{enumerate}

\begin{lstlisting}[language=Python]
hauteur = len(G)
largeur = len(G[0])
\end{lstlisting}

\begin{enumerate}
\def\labelenumi{\arabic{enumi}.}
\setcounter{enumi}{2}
\tightlist
\item
  \textbf{Réponse 3 :}
\end{enumerate}

\begin{lstlisting}[language=Python]
hauteur = len(G[0])
largeur = len(G[1])
\end{lstlisting}

\begin{enumerate}
\def\labelenumi{\arabic{enumi}.}
\setcounter{enumi}{3}
\tightlist
\item
  \textbf{Réponse 4 :}
\end{enumerate}

\begin{lstlisting}[language=Python]
hauteur = len(G[1])
largeur = len(G[0])
\end{lstlisting}

\end{exercice}

\begin{exercice}{}

Quelle est la valeur de l'expression
\passthrough{\lstinline![[0] * 3 for i in range(2)]!} ?

Réponses :

\begin{enumerate}
\def\labelenumi{\arabic{enumi}.}
\tightlist
\item
  \textbf{Réponse 1 :} \passthrough{\lstinline![[0,0], [0,0], [0,0]]!}
\item
  \textbf{Réponse 2 :} \passthrough{\lstinline![[0,0,0], [0,0,0]]!}
\item
  \textbf{Réponse 3 :} \passthrough{\lstinline![[0.000], [0.000]]!}
\item
  \textbf{Réponse 4 :}
  \passthrough{\lstinline![[0.00], [0.00], [0.00]]!}
\end{enumerate}

\end{exercice}

\begin{exercice}{}

On exécute le script suivant :

\begin{lstlisting}[language=Python]
a = [1, 2, 3]
b = [4, 5, 6]
c = a + b
\end{lstlisting}

Que contient la variable \passthrough{\lstinline!c!} à la fin de cette
exécution ?

Réponses :

\begin{enumerate}
\def\labelenumi{\arabic{enumi}.}
\item
  \textbf{Réponse 1 :} \passthrough{\lstinline![5,7,9]!}
\item
  \textbf{Réponse 2 :} \passthrough{\lstinline![1,4,2,5,3,6]!}
\item
  \textbf{Réponse 3 :} \passthrough{\lstinline![1,2,3,4,5,6]!}
\item
  \textbf{Réponse 4 :} \passthrough{\lstinline![1,2,3,5,7,9]!}
\end{enumerate}

\end{exercice}

\begin{exercice}{}

On exécute le script suivant~:

\begin{lstlisting}[language=Python]
asso = []
L=[['marc','marie'],['marie','jean'],
    ['paul','marie'], ['marie','marie'],
    ['marc','anne']]
for c in L :
    if c[1]=='marie':
        asso.append(c[0])
\end{lstlisting}

Que vaut asso à la fin de l'exécution ?

Réponses :

\begin{enumerate}
\def\labelenumi{\arabic{enumi}.}
\item
  \textbf{Réponses 1:}
  \passthrough{\lstinline!['marc', 'jean', 'paul']!}
\item
  \textbf{Réponses 2 :}
  \passthrough{\lstinline![['marc','marie'], ['paul','marie'], ['marie','marie']]!}
\item
  \textbf{Réponses 3 :}
  \passthrough{\lstinline!['marc', 'paul', 'marie']!}
\item
  \textbf{Réponses 4 :} \passthrough{\lstinline!['marie', 'anne']!}
\end{enumerate}

\end{exercice}

\begin{exercice}{}

Quelle est la valeur de la variable n à la fin de l'exécution du script
ci-dessous ?

\begin{lstlisting}[language=Python]
n = 1
while n != 20:
    n = n + 2
\end{lstlisting}

Réponses :

\begin{enumerate}
\def\labelenumi{\arabic{enumi}.}
\item
  \textbf{Réponse 1:} 1
\item
  \textbf{Réponse 2:} 20
\item
  \textbf{Réponse 3:} 22
\item
  \textbf{Réponse 4:} le programme ne termine pas, la boucle tourne
  indéfiniment
\end{enumerate}

\end{exercice}

\begin{exercice}{}

On définit en Python la fonction suivante :

\begin{lstlisting}[language=Python]
def f(L):
    S = []
    for i in range(len(L)-1):
        S.append(L[i] + L[i+1])
    return S
\end{lstlisting}

Quelle est la liste renvoyée par
\passthrough{\lstinline!f([1, 2, 3, 4, 5, 6])!}~?

\begin{enumerate}
\def\labelenumi{\arabic{enumi}.}
\tightlist
\item
  \passthrough{\lstinline![3, 5, 7, 9, 11, 13]!}
\item
  \passthrough{\lstinline![1, 3, 5, 7, 9, 11]!}
\item
  \passthrough{\lstinline![3, 5, 7, 9, 11]!}
\item
  cet appel de fonction déclenche un message d'erreur
\end{enumerate}

\end{exercice}

\hypertarget{repruxe9sentation-des-nombres-des-textes-ou-des-images}{%
\section{Représentation des nombres, des textes ou des
images}\label{repruxe9sentation-des-nombres-des-textes-ou-des-images}}

\begin{exercice}{}

Le codage en base deux de l'entier 26 en base dix est :

\begin{enumerate}
\def\labelenumi{\arabic{enumi}.}
\tightlist
\item
  11010
\item
  10010
\item
  11001
\item
  110010
\end{enumerate}

\end{exercice}

\begin{exercice}{}

Le résultat de la somme
\(\overline{101101}^{2} + \overline{101111}^{2}\) est :

\begin{enumerate}
\def\labelenumi{\arabic{enumi}.}
\tightlist
\item
  \(\overline{1100100}^{2}\)
\item
  \(\overline{1110101}^{2}\)
\item
  \(\overline{1011100}^{2}\)
\item
  \(\overline{1111100}^{2}\)
\end{enumerate}

\end{exercice}

\begin{exercice}{}

Le plus grand entier qu'on peut représenter en base deux sur 8 bits a
pour écriture décimale :

\begin{enumerate}
\def\labelenumi{\arabic{enumi}.}
\tightlist
\item
  11111111
\item
  10000000
\item
  255
\item
  256
\end{enumerate}

\end{exercice}

\begin{exercice}{}

Quelle est l'écriture décimale de l'entier qui s'écrit 1010 en binaire ?

Réponses :

\begin{enumerate}
\def\labelenumi{\arabic{enumi}.}
\item
  5
\item
  10
\item
  20
\item
  22
\end{enumerate}

\end{exercice}

\begin{exercice}{}

Combien d'entiers positifs ou nuls (entiers non signés) peut-on
représenter en machine sur 32~bits~?

Réponses :

\begin{enumerate}
\def\labelenumi{\arabic{enumi}.}
\tightlist
\item
  \(2^{32} - 1\)
\item
  \(2^{32}\)
\item
  \(2 \times 32\)
\item
  \(32^{2}\)
\end{enumerate}

\end{exercice}

\begin{exercice}{}

Combien de bits faut-il au minimum pour coder le nombre décimal 4085 ?

Réponses :

\begin{enumerate}
\def\labelenumi{\arabic{enumi}.}
\tightlist
\item
  4
\item
  12
\item
  2042
\item
  2043
\end{enumerate}

\end{exercice}

\begin{exercice}{}

Parmi les propositions suivantes, laquelle est la représentation binaire
de 761 ?

Réponses :

\begin{enumerate}
\def\labelenumi{\arabic{enumi}.}
\tightlist
\item
  11 1100 1101
\item
  11 1110 0101
\item
  10 0111 1001
\item
  10 1111 0001
\end{enumerate}

\end{exercice}

\begin{exercice}{}

Le codage d'une couleur se fait à l'aide de trois nombres compris
chacun, en écriture décimale, entre 0 et 255 (code RVB).

La couleur «~vert impérial~»~est codée, en écriture décimale, par (0,86,
27).

Le codage hexadécimal (base 16) correspondant est~:

Réponses :

\begin{enumerate}
\def\labelenumi{\arabic{enumi}.}
\tightlist
\item
  (0, 134, 39)
\item
  (0, 134, 1B)
\item
  (0, 56, 1B)
\item
  (0, 56, 39)
\end{enumerate}

\end{exercice}

\begin{exercice}{}

Quelle est l'écriture binaire du nombre entier 183 ?

Réponses :

\begin{enumerate}
\def\labelenumi{\arabic{enumi}.}
\tightlist
\item
  0100 1000
\item
  1110 1101
\item
  1011 0111
\item
  1001 0101
\end{enumerate}

\end{exercice}

\begin{exercice}{}

Le codage en base seize de l'entier \(256\) est :

\begin{enumerate}
\def\labelenumi{\arabic{enumi}.}
\tightlist
\item
  \(\overline{F1}^{16}\)
\item
  \(\overline{A1}^{16}\)
\item
  \(\overline{FA}^{16}\)
\item
  \(\overline{FF}^{16}\)
\end{enumerate}

\end{exercice}

\begin{exercice}{}

L'adresse MAC de la carte Wifi d'un smartphone est
\texttt{c8:60:00:a4:89:ab} avec six octets codés en base seize.

La transcription en base dix de cette adresse MAC est :

\begin{enumerate}
\def\labelenumi{\arabic{enumi}.}
\tightlist
\item
  \(200:96:0:164:137:171\)
\item
  \(20:6:0:14:17:21\)
\item
  \(96:0:0:40:72:110\)
\item
  \(140:6:0:74:152:186\)
\end{enumerate}

\end{exercice}

\begin{exercice}{}

La fonction ci-dessous doit retourner la liste des chiffres en base deux
d'un entier n par ordre décroissant des poids de gauche à droite.

\begin{lstlisting}[language=Python]
def nombre2chiffres(n):
    t = []
    while n >= 2:
        ..........
        n = n // 2
    ..............
    t.reverse()
    return t
\end{lstlisting}

Quelle instruction peut-on écrire en lignes 4 et 6 ?

\begin{enumerate}
\def\labelenumi{\arabic{enumi}.}
\tightlist
\item
  \passthrough{\lstinline!t.append(n)!}
\item
  \passthrough{\lstinline!t.append(n \% 2)!}
\item
  \passthrough{\lstinline!t.append(n // 2)!}
\item
  \passthrough{\lstinline!t = t + [n \% 2]!}
\end{enumerate}

\end{exercice}

\begin{exercice}{}

\begin{enumerate}
\def\labelenumi{\arabic{enumi}.}
\tightlist
\item
  Représenter en binaire le nombre d'écriture décimale 49.
\item
  Représenter en base dix, le nombre dont l'écriture en base deux est
  \(1010110\).
\item
  Déterminer le successeur en base deux des entiers :

  \begin{itemize}
  \tightlist
  \item
    \(111\)
  \item
    \(10011\)
  \item
    \(10111\)
  \end{itemize}
\item
  Déterminer le nombre de caractères qu'on peut coder sur un octet.
\end{enumerate}

\end{exercice}

\begin{exercice}{}

On souhaite convertir 25 de base 10 en base 2. On obtient en binaire :

Réponses :

\begin{enumerate}
\def\labelenumi{\arabic{enumi}.}
\tightlist
\item
  11001
\item
  10110
\item
  10011
\item
  11000
\end{enumerate}

\end{exercice}

\begin{exercice}{}

Quelle est la valeur affichée par l'exécution du test suivant ?

\begin{lstlisting}[language=Python]
In [1]: 0.1 + 0.2 == 0.3              
                                                                                                                                                                         
Out[1]: False
\end{lstlisting}

Réponses :

\begin{enumerate}
\def\labelenumi{\arabic{enumi}.}
\tightlist
\item
  \passthrough{\lstinline!True!}
\item
  \passthrough{\lstinline!False!}
\item
  0.3
\end{enumerate}

\end{exercice}

\begin{exercice}{}

Combien de bits sont nécessaires pour représenter 15 en binaire ?

Réponses :

\begin{enumerate}
\def\labelenumi{\arabic{enumi}.}
\tightlist
\item
  2
\item
  3
\item
  4
\item
  5
\end{enumerate}

\end{exercice}

\begin{exercice}{}

Quelle est la valeur affichée a l'exécution du programme Python suivant
?

\begin{lstlisting}[language=Python]
x = 1 
for i in range(11): 
    x = x * 2 
print(x)
\end{lstlisting}

Réponses :

\begin{enumerate}
\def\labelenumi{\arabic{enumi}.}
\tightlist
\item
  1024
\item
  2048
\item
  23
\item
  \(2^{2^{11}}\)
\end{enumerate}

\end{exercice}

\begin{exercice}{}

Quelle est la valeur affichée a l'exécution du programme Python suivant
?

\begin{lstlisting}[language=Python]
x = 2 
for i in range(10): 
    x = x ** 2 
print(x)
\end{lstlisting}

Réponses :

\begin{enumerate}
\def\labelenumi{\arabic{enumi}.}
\tightlist
\item
  1024
\item
  2048
\item
  \(2^{2^{10}}\)
\item
  \(2^{2^{11}}\)
\end{enumerate}

\end{exercice}

\begin{exercice}{}

\begin{enumerate}
\def\labelenumi{\arabic{enumi}.}
\tightlist
\item
  Convertir 3970 en base 6.
\item
  Convertir en base dix l'entier naturel dont l'écriture en base six est
  4321.
\item
  Écrire en Python une fonction \passthrough{\lstinline!base6(L)!} qui
  renvoie la valeur entière correspondant à la liste des chiffres de
  l'écriture en base 6.
\end{enumerate}

Exemple : \passthrough{\lstinline!base6([1, 3, 2])!} doit renvoyer 56
car \(1 \times 6^{2} + 3 \times 6 + 2 = 56\).

\end{exercice}

\begin{exercice}{}

\begin{enumerate}
\def\labelenumi{\arabic{enumi}.}
\tightlist
\item
  Pour déterminer la liste des chiffres en base dix d'un entier naturel,
  un élève a écrit la fonction ci-dessous :
\end{enumerate}

\begin{lstlisting}[language=Python]
def liste_chiffres(n):
    L = [n % 10]
    while n > 0:
        n = n // 10
        L.insert(0, n % 10)
    return L
\end{lstlisting}

Malheureusement sa fonction ne retourne pas le résultat attendu pour
l'entier \(730\) :

\begin{lstlisting}[language=Python]
>>> liste_chiffres(730)
[0, 7, 3, 0]
\end{lstlisting}

Proposer une version corrigée de la fonction
\passthrough{\lstinline!liste\\\_chiffres!}.

\begin{enumerate}
\def\labelenumi{\arabic{enumi}.}
\setcounter{enumi}{1}
\tightlist
\item
  Compléter la fonction
  \passthrough{\lstinline!somme\_chiffres\_base2(n)!} pour qu'elle
  retourne la somme des chiffres en base deux de l'entier
  \passthrough{\lstinline!n!} passé en paramètre.
\end{enumerate}

\begin{lstlisting}[language=Python]
def somme_chiffres_base2(n):
    s = 0
    while n > 0:
        s = s + n % 2
        ..............
    return s
\end{lstlisting}

\begin{enumerate}
\def\labelenumi{\arabic{enumi}.}
\setcounter{enumi}{2}
\tightlist
\item
  Déterminer une valeur possible de la variable \texttt{secret} telle
  que :
\end{enumerate}

\begin{lstlisting}[language=Python]
>>> somme_chiffres(secret)
734
\end{lstlisting}

\begin{enumerate}
\def\labelenumi{\arabic{enumi}.}
\setcounter{enumi}{3}
\tightlist
\item
  Écrire une fonction \passthrough{\lstinline!maximum\_chiffre(n)!} qui
  retourne le plus grand chiffre de l'écriture en base dix de l'entier
  naturel \passthrough{\lstinline!n!} passé en paramètre.
\end{enumerate}

\end{exercice}

\begin{exercice}{}

Le nombre de chiffres en base \(2\) d'un entier naturel \(n\) est :

\begin{itemize}
\tightlist
\item
  \(1\) si \(n\) est égal à \(0\) ;
\item
  l'unique entier \(p\) tel que \(2^{p-1} \leqslant n < 2^{p}\) si \(n\)
  supérieur à \(0\).
\end{itemize}

Compléter la fonction ci-dessous pour qu'elle retourne le nombre de
chiffres en base \(2\) de l'entier \texttt{n} passé en paramètres.

\begin{lstlisting}[language=Python]
def nbchiffres_base2(n):
    if n == 0:
        return 1
    nbchiffre = 1
    puissance = 1
    while ..............:
        .........................
        .........................
    return ......................
\end{lstlisting}

\end{exercice}

\hypertarget{logique-booluxe9enne-et-architecture}{%
\section{Logique booléenne et
architecture}\label{logique-booluxe9enne-et-architecture}}

\begin{exercice}{}

En 1945, John Von Neumann a décrit un modèle d'architecture qui est
encore valable pour les ordinateurs actuels. Quelles sont les entités de
ce modèle et comment communiquent-elles ?

\end{exercice}

\begin{exercice}{}

Parmi les quatre expressions suivantes, laquelle s'évalue en True ?

Réponses :

\begin{enumerate}
\def\labelenumi{\arabic{enumi}.}
\item
  \textbf{Réponse 1 :}
  \passthrough{\lstinline!False and (True and False)!}
\item
  \textbf{Réponse 2 :}
  \passthrough{\lstinline!False or (True and False)!}
\item
  \textbf{Réponse 3 :}
  \passthrough{\lstinline!True and (True and False)!}
\item
  \textbf{Réponse 4 :}
  \passthrough{\lstinline!True or (True and False)!}
\end{enumerate}

\end{exercice}

\begin{exercice}{}

Sachant que l'expression \passthrough{\lstinline!not(a or b)!} a la
valeur True, quelles peuvent être les valeurs des variables booléennes a
et b~?

Réponses :

\begin{enumerate}
\def\labelenumi{\arabic{enumi}.}
\tightlist
\item
  True et True
\item
  False et True
\item
  True et False
\item
  False et False
\end{enumerate}

\end{exercice}

\begin{exercice}{}

A quelle expression logique correspond cette table de vérité ?

\begin{longtable}[]{@{}cll@{}}
\toprule
A & B & f(A, B)\tabularnewline
\midrule
\endhead
0 & 0 & 1\tabularnewline
0 & 1 & 1\tabularnewline
1 & 0 & 1\tabularnewline
1 & 1 & 0\tabularnewline
\bottomrule
\end{longtable}

\end{exercice}

\begin{exercice}{}

A quelle expression logique correspond cette table de vérité ?

\begin{longtable}[]{@{}cll@{}}
\toprule
A & B & f(A, B)\tabularnewline
\midrule
\endhead
0 & 0 & 0\tabularnewline
0 & 1 & 0\tabularnewline
1 & 0 & 0\tabularnewline
1 & 1 & 1\tabularnewline
\bottomrule
\end{longtable}

\end{exercice}

\begin{exercice}{}

A quelle expression logique correspond cette table de vérité ?

\begin{longtable}[]{@{}cll@{}}
\toprule
A & B & f(A, B)\tabularnewline
\midrule
\endhead
0 & 0 & 0\tabularnewline
0 & 1 & 1\tabularnewline
1 & 0 & 1\tabularnewline
1 & 1 & 1\tabularnewline
\bottomrule
\end{longtable}

\end{exercice}

\begin{exercice}{}

A quelle expression logique correspond cette table de vérité ?

\begin{longtable}[]{@{}cll@{}}
\toprule
A & B & f(A, B)\tabularnewline
\midrule
\endhead
0 & 0 & 0\tabularnewline
0 & 1 & 1\tabularnewline
1 & 0 & 1\tabularnewline
1 & 1 & 0\tabularnewline
\bottomrule
\end{longtable}

\end{exercice}

\begin{exercice}{}

Choisir une expression booléenne pour la variable S qui satisfait la
table de vérité suivante.

\begin{longtable}[]{@{}lll@{}}
\toprule
A & B & S\tabularnewline
\midrule
\endhead
0 & 0 & 1\tabularnewline
0 & 1 & 0\tabularnewline
1 & 0 & 1\tabularnewline
1 & 1 & 1\tabularnewline
\bottomrule
\end{longtable}

Réponse :

\begin{enumerate}
\def\labelenumi{\arabic{enumi}.}
\tightlist
\item
  A ou (non B)
\item
  (non A) ou B
\item
  (non A) ou (non B)
\item
  non (A ou B)
\end{enumerate}

\end{exercice}

\hypertarget{algorithmique}{%
\section{Algorithmique}\label{algorithmique}}

\begin{exercice}{}

T est un tableau de nombres entiers non vide. Que représente la valeur
de s~renvoyée par cette fonction ?

\begin{lstlisting}[language=Python]
def mystere(T):
    s = 0
    for k in T:
        if k % 2 == 0:
            s = s+k
    return s
\end{lstlisting}

Réponses :

\begin{enumerate}
\def\labelenumi{\arabic{enumi}.}
\item
  \textbf{Réponse 1 :} la somme des valeurs du tableau T
\item
  \textbf{Réponse 2 :} la somme des valeurs positives du tableau T
\item
  \textbf{Réponse 3 :} la somme des valeurs impaires du tableau T
\item
  \textbf{Réponse 4 :} la somme des valeurs paires du tableau T
\end{enumerate}

\end{exercice}

\begin{exercice}{}

Soit \(n\) un entier naturel. Sa factorielle est le produit des nombres
entiers strictement positifs qui sont plus petits ou égaux à \(n\). Par
exemple la factorielle de 4 vaut \(1 \times 2 \times 3 \times 4 = 24\).

Quelle est la fonction correcte parmi les suivantes ?

Réponses :

\begin{enumerate}
\def\labelenumi{\arabic{enumi}.}
\tightlist
\item
  \textbf{Réponse 1 :}
\end{enumerate}

\begin{lstlisting}[language=Python]
def factorielle(n):
    i = 0
    fact = 1
    while i <= n:
        fact = fact * i
        i = i + 1
    return fact
\end{lstlisting}

\begin{enumerate}
\def\labelenumi{\arabic{enumi}.}
\setcounter{enumi}{1}
\tightlist
\item
  \textbf{Réponse 2 :}
\end{enumerate}

\begin{lstlisting}[language=Python]
def factorielle(n):
    i = 1
    fact = 1
    while i < n:
        fact = fact * i
        i = i + 1
    return fact
\end{lstlisting}

\begin{enumerate}
\def\labelenumi{\arabic{enumi}.}
\setcounter{enumi}{2}
\tightlist
\item
  \textbf{Réponse 3 :}\\
\end{enumerate}

\begin{lstlisting}[language=Python]
def factorielle(n):
    i = 0
    fact = 1
    while i < n:
        i = i + 1
        fact = fact * i
    return fact
\end{lstlisting}

\begin{enumerate}
\def\labelenumi{\arabic{enumi}.}
\setcounter{enumi}{3}
\tightlist
\item
  \textbf{Réponse 4 :}
\end{enumerate}

\begin{lstlisting}[language=Python]
def factorielle(n):
    i = 0
    fact = 1
    while i <= n:
        i = i + 1
        fact = fact * i
    return fact
\end{lstlisting}

\end{exercice}

\begin{exercice}{}

Écrire une fonction \passthrough{\lstinline!min2(a, b)!} qui retourne le
minimum de deux nombres passés en paramètre, sans utiliser la fonction
min du module builtins.

\end{exercice}

\begin{exercice}{}

En voulant programmer une fonction qui calcule la valeur minimale d'une
liste d'entiers, on a écrit :

\begin{lstlisting}[language=Python]
def minimum(L):
    mini = 0
    for e in L:
        if e < mini:
            mini = e
    return mini
\end{lstlisting}

Cette fonction a été mal programmée. Pour quelle liste ne donnera-t-elle
pas le résultat attendu, c'est-à-dire son minimum ?

Réponses :

\begin{enumerate}
\def\labelenumi{\arabic{enumi}.}
\tightlist
\item
  \textbf{Réponse 1 :} \passthrough{\lstinline![-1,-8,12,2,23]!}
\item
  \textbf{Réponse 2:} \passthrough{\lstinline![0,18,12,2,3]!}
\item
  \textbf{Réponse 3:} \passthrough{\lstinline![-1,-1,12,12,23]!}
\item
  \textbf{Réponse 4:} \passthrough{\lstinline![1,8,12,2,23]!}
\end{enumerate}

\end{exercice}

\begin{exercice}{}

Écrire une fonction \passthrough{\lstinline!min\_liste(L)!} qui retourne
le minimum d'une liste de nombres passée en paramètre.

\end{exercice}

\begin{exercice}{}

Quelle est la valeur de element à la fin de l'exécution du code suivant
:

\begin{lstlisting}[language=Python]
L = [1,2,3,4,1,2,3,4,0,2]
element = L[0]
for k in L:
    if k > element:
        element = k
\end{lstlisting}

\begin{enumerate}
\def\labelenumi{\arabic{enumi}.}
\tightlist
\item
  \textbf{Réponse 1 :} 0
\item
  \textbf{Réponse 2:} 1
\item
  \textbf{Réponse 3:} 4
\item
  \textbf{Réponse 4:} 10
\end{enumerate}

\end{exercice}

\begin{exercice}{}

Un élève a écrit la fonction ci-dessus pour déterminer si le premier
paramètre n est dans la liste L passée en second paramètre. Le
professeur lui indique que son code comporte des erreurs. Proposer une
version corrigée de cette fonction.

\begin{lstlisting}[language=Python]
def element_dans_liste(n, L):
    for k in range(len(L)):
        if e == n:
            return True
        else:
            return False
\end{lstlisting}

\end{exercice}

\begin{exercice}{}

Pour chacune des listes ci-dessous, déterminer si l'algorithme le plus
adapté pour qu'une machine y recherche une valeur est l'algorithme de
recherche séquentielle ou l'algorithme de recherche dichotomique.

\begin{enumerate}
\def\labelenumi{\arabic{enumi}.}
\tightlist
\item
  Liste A :
  \passthrough{\lstinline![4, 6, 15, 20, 21, 26, 31, 41, 42, 50, 69, 87, 88, 92, 97]!}
\item
  Liste B :
  \passthrough{\lstinline![41, 97, 91, 59, 7, 45, 3, 96, 26, 32, 18, 11, 67, 74, 54]!}
\end{enumerate}

\end{exercice}

\begin{exercice}{}

Compléter le code de la fonction
\passthrough{\lstinline!recherche\_dicho\_dec(x, L)!} qui prend en
paramètres un nombre \passthrough{\lstinline!x!} et une liste de nombres
\passthrough{\lstinline!L!} triée dans l'ordre décroissant. Elle doit
retourner \passthrough{\lstinline!True!} si \passthrough{\lstinline!x!}
appartient à L et False sinon.

\begin{lstlisting}[language=Python]
def recherche_dicho_dec(x, L):
    a, b = 0, len(L) - 1
    while a <= b:
        m = (a + b) // 2
        if L[m] > x:
            ..................
        elif L[m] < x:
            ................
        else:
            .................
    .........................
\end{lstlisting}

\end{exercice}

\begin{exercice}{}

\emph{Auteur Sylvie Genre}

La fonction suivante doit calculer la moyenne d'un tableau de nombres,
passé en paramètre. Avec quelles expressions faut-il remplacer les
points de suspension pour que la fonction soit correcte ?

\begin{lstlisting}[language=Python]
def moyenne(tableau):
    total = ...
    for valeur in tableau:
        total = total + valeur
    return total / ...
\end{lstlisting}

Réponses :

\begin{enumerate}
\def\labelenumi{\arabic{enumi}.}
\tightlist
\item
  \textbf{Réponse 1 :} 1 et \passthrough{\lstinline!(len(tableau) + 1)!}
\item
  \textbf{Réponse 2:} 0 et \passthrough{\lstinline!(len(tableau) + 1)!}
\item
  \textbf{Réponse 3:} 0 et \passthrough{\lstinline!len(tableau)!}
\end{enumerate}

\end{exercice}

\begin{exercice}{}

On considère le code incomplet suivant qui recherche le maximum dans une
liste.

\begin{lstlisting}[language=Python]
liste = [5,12,15,3,15,17,29,1]
iMax = 0
for i in range(1,len(liste)):
    .........
iMax = i
print (liste[iMax])
\end{lstlisting}

Par quoi faut-il remplacer la ligne pointillée ?

Réponses :

\begin{enumerate}
\def\labelenumi{\arabic{enumi}.}
\tightlist
\item
  \textbf{Réponse 1 :} \passthrough{\lstinline!if i > iMax:!}
\item
  \textbf{Réponse 2:}
  \passthrough{\lstinline!if liste[i] > liste[iMax]:!}
\item
  \textbf{Réponse 3:} \passthrough{\lstinline!if liste[i] > iMax:!}
\item
  \textbf{Réponse 4:} \passthrough{\lstinline!if i > liste[iMax]:!}
\end{enumerate}

\end{exercice}

\hypertarget{le-web}{%
\section{Le Web}\label{le-web}}

\begin{exercice}{}

Le protocole qui permet à un client de faire une requête de page Web
auprès d'un serveur Web s'appelle :

Réponses :

\begin{enumerate}
\def\labelenumi{\arabic{enumi}.}
\tightlist
\item
  Internet Protocol
\item
  HTML
\item
  HTTP
\item
  TCP
\item
  WWW
\end{enumerate}

\end{exercice}

\begin{exercice}{}

Si la page demandée n'est pas disponible, le serveur Web renvoie au
client un code d'erreur :

Réponses :

\begin{enumerate}
\def\labelenumi{\arabic{enumi}.}
\item
  404
\item
  504
\item
  403
\item
  503
\end{enumerate}

\end{exercice}

\begin{exercice}{}

L'inventeur du World Wide Web au CERN est :

Réponses :

\begin{enumerate}
\def\labelenumi{\arabic{enumi}.}
\item
  Tim Berners-Lee
\item
  Ted Nelson
\item
  Louis Pouzin
\item
  Vinton Cerf
\end{enumerate}

\end{exercice}

\begin{exercice}{}

Dans le fichier source d'une page Web, le code qui permet de créer un
lien hypertexte vers la page \url{https://www.w3schools.com/} est :

Réponses :

\begin{enumerate}
\def\labelenumi{\arabic{enumi}.}
\item
  \passthrough{\lstinline!<a href="https://www.w3schools.com/">lien hypertexte</a>!}
\item
  \passthrough{\lstinline!<a href="https://www.w3schools.com/">https://www.w3schools.com/</a>!}
\item
  \passthrough{\lstinline!<a href="lien hypertexte">https://www.w3schools.com/</a>!}
\item
  \passthrough{\lstinline!<href a="https://www.w3schools.com/">lien hypertexte</href>!}
\end{enumerate}

\end{exercice}

\begin{exercice}{}

Pour créer un lien vers la page d'accueil de Wikipédia, que devra-t-on
écrire dans une page Web ?

\textbf{Réponses :}

\begin{enumerate}
\def\labelenumi{\arabic{enumi}.}
\item
  \textbf{Réponse 1 :}
  \passthrough{\lstinline!<a target="http://fr.wikipedia.org">Wikipédia</a>!}
\item
  \textbf{Réponse 2 :}
  \passthrough{\lstinline!<a href="http://fr.wikipedia.org">!}
\item
  \textbf{Réponse 3 :}
  \passthrough{\lstinline!<a href="http://fr.wikipedia.org">Wikipédia</a>!}
\item
  \textbf{Réponse 4 :}
  \passthrough{\lstinline!<link src="http://fr.wikipedia.org">Wikipédia</link>!}
\end{enumerate}

\end{exercice}

\end{document}
